\Article{39}
To make clearer how far it is legitimate to equate the value of \(M\) arising by chance with its mean, it is still necessary to investigate the mean error to be feared when we put \(\frac{M}{\pi-\rho}=\mu\mu\). This mean error is equal to the square root of the mean value of the quantity
\[
\left(\frac{\Omega^0-x^0\xi^0-y^0\eta^0-z^0\zeta^0-\text{etc.}-(\pi-\rho)\mu\mu}{\pi-\rho}\right)^2
\]
which we shall express as follows:
\[
\begin{aligned}
&\left(\frac{\Omega^0-x^0\xi^0-y^0\eta^0-z^0\zeta^0-\text{etc.}}{\pi-\rho}\right)^2\\[4pt]
&\quad-\frac{2\mu\mu}{\pi-\rho}(\Omega^0-x^0\xi^0-y^0\eta^0-z^0\zeta^0-\text{etc.}-(\pi-\rho)\mu\mu)-\mu^4
\end{aligned}
\]
and since the mean value of the second term is clearly \(=0\), what we have to investigate is the mean value of the function
\[
\Psi=(\Omega^0-x^0\xi^0-y^0\eta^0-z^0\zeta^0-\text{etc.})^2
\]
Once this has been found and denoted by \(N\), the required mean error will be
\[
=\BareRoot{\left(\frac{N}{(\pi-\rho)^2}-\mu^4\right)}
\]
The expanded expression \(\Psi\) is clearly a homogeneous function either of the errors \(e,e',e''\) etc. or of the quantities \(\varepsilon,\varepsilon',\varepsilon''\) etc., and its mean value will be found if

\(1^\mathrm{o}\) the fourth powers \(e^4,e'^4,e''^4\) etc. are replaced by their mean values

\(2^\mathrm{o}\) each product of two squares, such as \(eee'e',\allowbreak eee''e'',\allowbreak e'e'e''e''\) etc., is replaced by the product of their mean values, namely \(mmm'm',\allowbreak mmm''m'',\allowbreak m'm'm''m''\) etc.
\par\SourceLeaf{055}%
\(3^\mathrm{o}\) all the remaining parts, which will involve either a factor such as \(e^3e'\) or one such as \(eee'e''\), are omitted altogether.

We shall assume the mean values of the fourth powers \(e^4,\allowbreak e'^4,\allowbreak e''^4\) etc. to be proportional to the fourth powers \(m^4,\allowbreak m'^4,\allowbreak m''^4\) etc. themselves (see art.\ 16), so that the former are to the latter as \(\nu^4\) to \(\mu^4\); thus \(\nu^4\) denotes the mean value of the fourth powers of observations whose weight is \(=1\). The preceding rules can therefore also be expressed as follows: in place of each fourth power \(\varepsilon^4,\varepsilon'^4,\varepsilon''^4\) etc. we must write \(\nu^4\), in place of each product of two squares, such as \(\varepsilon\varepsilon\varepsilon'\varepsilon',\allowbreak\varepsilon\varepsilon\varepsilon''\varepsilon'',\allowbreak\varepsilon'\varepsilon'\varepsilon''\varepsilon''\) etc., we must write \(\mu^4\), and all remaining terms that will involve factors such as \(\varepsilon^3\varepsilon'\) or \(\varepsilon\varepsilon\varepsilon'\varepsilon''\) or \(\varepsilon\varepsilon'\varepsilon''\varepsilon'''\) must be suppressed.

With this properly understood, it will readily be seen that

I. The mean value of the square \(\Omega^0\Omega^0\) is \(\pi\nu^4+(\pi\pi-\pi)\mu^4\)

II. The mean value of the product \(\varepsilon\varepsilon x^0\xi^0\) is \(=a\alpha\nu^4+(a'\alpha'+a''\alpha''+\text{etc.})\mu^4\), or, since \(a\alpha+a'\alpha'+a''\alpha''+\text{etc.}=1\),
\[
=a\alpha(\nu^4-\mu^4)+\mu^4
\]
And since, similarly, the mean value of the product \(\varepsilon'\varepsilon'x^0\xi^0\) is \(=a'\alpha'(\nu^4-\mu^4)+\mu^4\), while the mean value of the product \(\varepsilon''\varepsilon''x^0\xi^0\) is \(=a''\alpha''(\nu^4-\mu^4)+\mu^4\) and so on, it is clear that the mean value of the product \((\varepsilon\varepsilon+\varepsilon'\varepsilon'+\varepsilon''\varepsilon''+\text{etc.})x^0\xi^0\) or \(\Omega^0x^0\xi^0\) is
\[
=\nu^4-\mu^4+\pi\mu^4
\]
\begingroup\setlength{\rightskip}{0pt plus 2em}\noindent The products \(\Omega^0y^0\eta^0,\allowbreak\Omega^0z^0\zeta^0\) etc. will have the same mean value. Therefore the mean value of the product \(\Omega^0(x^0\xi^0+y^0\eta^0+z^0\zeta^0+\text{etc.})\) is\par\endgroup
\[
=\rho\nu^4+\rho(\pi-1)\mu^4
\]

III. To keep the remaining expansions from becoming too lengthy, a suitable notation must be introduced. We shall therefore use the symbol \(\Sigma\) in a somewhat wider sense than has been done at various places above: it will denote the sum of the term to which it is prefixed and all similar but nonidentical terms arising from it by all permutations of the observations. In this way, e.\ g., we have \(x^0=\Sigma\alpha\varepsilon\), \(x^0x^0=\Sigma\alpha\alpha\varepsilon\varepsilon+2\Sigma\alpha\alpha'\varepsilon\varepsilon'\). Collecting the mean value of the product \(x^0x^0\xi^0\xi^0\) by parts, we therefore first have the mean value of the product \(\alpha\alpha\varepsilon\varepsilon\xi^0\xi^0\)
\[
\begin{aligned}
&=aa\alpha\alpha\nu^4+\alpha\alpha(a'a'+a''a''+\text{etc.})\mu^4\\
&=aa\alpha\alpha(\nu^4-\mu^4)+\alpha\alpha\mu^4\Sigma aa
\end{aligned}
\]
\SourceLeaf{056}
Similarly the mean value of the product \(\alpha'\alpha'\varepsilon'\varepsilon'\xi^0\xi^0\) is \(=a'a'\alpha'\alpha'(\nu^4-\mu^4)+\alpha'\alpha'\mu^4\Sigma aa\) and so on, and therefore the mean value of the product \(\xi^0\xi^0\Sigma\alpha\alpha\varepsilon\varepsilon\) is
\[
=(\nu^4-\mu^4)\Sigma aa\alpha\alpha+\mu^4\Sigma aa.\Sigma\alpha\alpha
\]
Further, the mean value of the product \(\alpha\alpha'\varepsilon\varepsilon'\xi^0\xi^0\) is \(=2aa'\alpha\alpha'\mu^4\), and similarly the mean value of the product \(\alpha\alpha''\varepsilon\varepsilon''\xi^0\xi^0\) is \(=2aa''\alpha\alpha''\mu^4\) etc.; from this it is readily concluded that the mean value of the product \(\xi^0\xi^0\Sigma\alpha\alpha'\varepsilon\varepsilon'\) is
\[
\begin{aligned}
&=2\mu^4\Sigma a\alpha a'\alpha'=\mu^4((\Sigma a\alpha)^2-\Sigma aa\alpha\alpha)\\
&=\mu^4(1-\Sigma aa\alpha\alpha)
\end{aligned}
\]
Collecting these, we have the mean value of the product \(x^0x^0\xi^0\xi^0\)
\[
=(\nu^4-3\mu^4)\Sigma aa\alpha\alpha+2\mu^4+\mu^4\Sigma aa.\Sigma\alpha\alpha
\]

IV. In a very similar way the mean value of the product \(x^0y^0\xi^0\eta^0\) is found to be
\[
\begin{aligned}
&=\nu^4\Sigma ab\alpha\beta+\mu^4\Sigma a\alpha b'\beta'\\
&\quad+\mu^4\Sigma ab\alpha'\beta'+\mu^4\Sigma a\beta b'\alpha'
\end{aligned}
\]
But the following holds
\[
\begin{aligned}
\Sigma a\alpha b'\beta'&=\Sigma a\alpha.\Sigma b\beta-\Sigma a\alpha b\beta\\
\Sigma ab\alpha'\beta'&=\Sigma ab.\Sigma\alpha\beta-\Sigma ab\alpha\beta\\
\Sigma a\beta b'\alpha'&=\Sigma a\beta.\Sigma b\alpha-\Sigma a\beta b\alpha
\end{aligned}
\]
therefore, because \(\Sigma a\alpha=1,\allowbreak\Sigma b\beta=1,\allowbreak\Sigma a\beta=0,\allowbreak\Sigma b\alpha=0\), that mean value becomes
\[
=(\nu^4-3\mu^4)\Sigma ab\alpha\beta+\mu^4(1+\Sigma ab.\Sigma\alpha\beta)
\]

V. Since, in exactly the same way, the mean value of the product \(x^0z^0\xi^0\zeta^0\) is
\[
=(\nu^4-3\mu^4)\Sigma ac\alpha\gamma+\mu^4(1+\Sigma ac.\Sigma\alpha\gamma)
\]
and so on, addition gives the mean value of the product \(x^0\xi^0(x^0\xi^0+y^0\eta^0+z^0\zeta^0+\text{etc.})\) as
\[
\begin{aligned}
&=(\nu^4-3\mu^4)\Sigma(a\alpha(a\alpha+b\beta+c\gamma+\text{etc.}))+(\rho+1)\mu^4\\
&\quad+\mu^4(\Sigma aa.\Sigma\alpha\alpha+\Sigma ab.\Sigma\alpha\beta+\Sigma ac.\Sigma\alpha\gamma+\text{etc.})\\
&=(\nu^4-3\mu^4)\Sigma(a\alpha(a\alpha+b\beta+c\gamma+\text{etc.}))+(\rho+2)\mu^4
\end{aligned}
\]

VI. In exactly the same way the mean value of the product \(y^0\eta^0(x^0\xi^0+y^0\eta^0+z^0\zeta^0+\text{etc.})\) is found to be
\SourceLeaf{057}
\[
=(\nu^4-3\mu^4)\Sigma(b\beta(a\alpha+b\beta+c\gamma+\text{etc.}))+(\rho+2)\mu^4
\]
then the mean value of the product \(z^0\zeta^0(x^0\xi^0+y^0\eta^0+z^0\zeta^0+\text{etc.})\) is
\[
=(\nu^4-3\mu^4)\Sigma(c\gamma(a\alpha+b\beta+c\gamma+\text{etc.}))+(\rho+2)\mu^4
\]
and so on. Addition therefore gives the mean value of the square \((x^0\xi^0+y^0\eta^0+z^0\zeta^0+\text{etc.})^2\)
\[
=(\nu^4-3\mu^4)\Sigma((a\alpha+b\beta+c\gamma+\text{etc.})^2)+(\rho\rho+2\rho)\mu^4
\]

VII. Finally, collecting everything properly gives
\[
\begin{aligned}
N&=(\pi-2\rho)\nu^4+(\pi\pi-\pi-2\pi\rho+4\rho+\rho\rho)\mu^4\\
&\quad+(\nu^4-3\mu^4)\Sigma((a\alpha+b\beta+c\gamma+\text{etc.})^2)\\
&=(\pi-\rho)(\nu^4-\mu^4)+(\pi-\rho)^2\mu^4\\
&\quad-(\nu^4-3\mu^4)[\rho-\Sigma((a\alpha+b\beta+c\gamma+\text{etc.})^2)]
\end{aligned}
\]
Thus the mean error to be feared in determining \(\mu\mu\) by the formula
\[
\mu\mu=\frac{M}{\pi-\rho}
\]
will be
\[
=\BareRoot{\left\{\frac{\nu^4-\mu^4}{\pi-\rho}-\frac{\nu^4-3\mu^4}{(\pi-\rho)^2}.[\rho-\Sigma((a\alpha+b\beta+c\gamma+\text{etc.})^2)]\right\}}
\]

\Article{40}
The quantity \(\Sigma((a\alpha+b\beta+c\gamma+\text{etc.})^2)\), which enters the expression just found, cannot in general be reduced to a simpler form. Nevertheless two bounds can be assigned between which its value must necessarily lie. \textit{First}, from the relations developed above it is readily proved that
\[
\begin{aligned}
&(a\alpha+b\beta+c\gamma+\text{etc.})^2+(a\alpha'+b\beta'+c\gamma'+\text{etc.})^2\\
&\quad+(a\alpha''+b\beta''+c\gamma''+\text{etc.})^2+\text{etc.}\\
&=a\alpha+b\beta+c\gamma+\text{etc.}
\end{aligned}
\]
from which we conclude that \(a\alpha+b\beta+c\gamma+\text{etc.}\) is a positive quantity smaller than unity (at least not greater). The same holds for the quantity \(a'\alpha'+b'\beta'+c'\gamma'+\text{etc.}\), which is found to be equal to the sum
\[
\begin{aligned}
&(a'\alpha+b'\beta+c'\gamma+\text{etc.})^2+(a'\alpha'+b'\beta'+c'\gamma'+\text{etc.})^2\\
&\quad+(a'\alpha''+b'\beta''+c'\gamma''\text{ etc.})^2+\text{etc.}
\end{aligned}
\]
and similarly \(a''\alpha''+b''\beta''+c''\gamma''+\text{etc.}\) will be smaller than unity, and so
\SourceLeaf{058}%
on. Hence \(\Sigma((a\alpha+b\beta+c\gamma+\text{etc.})^2)\) is necessarily smaller than \(\pi\). \textit{Second}, we have \(\Sigma(a\alpha+b\beta+c\gamma+\text{etc.})=\rho\), since \(\Sigma a\alpha=1,\allowbreak\Sigma b\beta=1,\allowbreak\Sigma c\gamma=1\) etc.; from this it is readily deduced that the sum of squares \(\Sigma((a\alpha+b\beta+c\gamma+\text{etc.})^2)\) is greater than \(\frac{\rho\rho}{\pi}\), or at least not smaller. Thus the term
\[
\frac{\nu^4-3\mu^4}{(\pi-\rho)^2}.[\rho-\Sigma((a\alpha+b\beta+c\gamma+\text{etc.})^2)]
\]
necessarily lies between the bounds \(-\frac{\nu^4-3\mu^4}{\pi-\rho}\) and \(\frac{\nu^4-3\mu^4}{\pi-\rho}.\frac{\rho}{\pi}\) or, if we prefer wider ones, between \(-\frac{\nu^4-3\mu^4}{\pi-\rho}\) and \(+\frac{\nu^4-3\mu^4}{\pi-\rho}\), and therefore the square of the mean error to be feared in the value of \(\mu\mu=\frac{M}{\pi-\rho}\) lies between the bounds \(\frac{2\nu^4-4\mu^4}{\pi-\rho}\) and \(\frac{2\mu^4}{\pi-\rho}\), so that any degree of precision can be attained, provided only that the number of observations is sufficiently large.

It is very remarkable that under the hypothesis (art.\ 9, III) on which the theory of least squares was formerly based, that term disappears altogether; and just as, to obtain an approximate value of the mean error of the observations \(\mu\), in all cases the sum \(\lambda\lambda+\lambda'\lambda'+\lambda''\lambda''+\text{etc.}=M\) must be treated as if it were a sum of \(\pi-\rho\) errors arising by chance, so under that hypothesis the precision of this determination itself also becomes equal to that which we found should be assigned to a determination from \(\pi-\rho\) true errors in art.\ 15.
\par\nobreak\medskip\centerline{\rule{.25\linewidth}{.3pt}}
\clearpage
